% !TeX root = ../Western Philosophy.tex
\section{笛卡尔哲学}

\subsection{笛卡尔生平简介}

笛卡尔幼年体弱多病（真是毫无意思的名人开头），24-29岁周游列国。
31岁以后转向数学物理，有正弦定律，解析几何和光学的工作

3，40年代笛卡尔转向哲学，在《Of the World》出版作罢之后，出版著作
《谈谈方法》(1637)，序文表达了基础的哲学观点。自称为体系性的哲学，体系的基础是
形而上学（以及数学），体系依次向上，有物理学等实践科学和道德哲学等。

1641年出版\textbf{《第一哲学沉思》}（形而上学），48岁出版《哲学原理》（物理），
去世前作品《论灵魂的激情》（道德哲学）注意哲学上的的激情指的是“受动者”，和主动（Action）相对

\subsection{Discart的我思}

我思提出的目的是应对怀疑论的挑战，为科学认识提供基础（Foundationalism）。笛卡尔提出“我思故我在”。论证的逻辑是即使有Demon在欺骗我，只要我在思考的时候，我就存在。

\textbf{Cogito，ergo sum.（I think， therefore I am）}

\begin{fancybox}{文本}
    Undoubtedly I exist if he deceives me; let him deceive me as much as he can,he will never bring it about that I am nothing while I'm thinking that I'm something. 
\end{fancybox}

我思指的不是Discart本人，而是心灵与诸多心灵状态。

\begin{itemize}
    \item 我思是一个思维的存在着
    \item 心灵被重新界定：意识成为了心灵的确定性标准。
    
    需要注意，具有意识在之前是功能性的标准：具备感官，有理智的心智。能够完成这些功能（感知，推理）的东西就具有心灵。在Discart看来，动物是没有心灵的，理由是动物没有意识，是一个活动的机械体（？）
    \item 心灵状态的含义非常广泛，包括思考，内省等等。
\end{itemize}

一些具体的心灵状态论述：

\begin{itemize}
    \item 判断属于意志类而非理智类活动，4边从而解释了错误的出现
    \item 想象和理智不同（考虑：三边形，五边形，千边形，万边形）
    \item \textbf{感觉可以独立于身体对象和外部世界存在（只要求心灵）}
    
    既然有了Demon一类东西的存在，则外部世界未必存在（什么都没有），人们的感觉都是对思维状态和心灵状态的感知。所谓对外部的认识是从心灵状态外推的。
\end{itemize}

\subsection{上帝}

Discart在Meditation中出现了多处关于上帝的论证，上帝的引入是必要的：由于前面的论证只是说明了心灵状态idea和thought的存在，无法解决外部世界怀疑论的问题。Discart为了证明外部世界存在，构想出了上帝。他将观念分为三类：

\begin{itemize}
    \item 天赋的：thing，truth等观念是先天具备的。
    \item 获得的
    \item 自创的
\end{itemize}

对于上帝的证明有以下几种：

第一种，Discart认为上帝的观念是获得的，因为上帝的观念是完满的，因此一定来自于上帝。《第二沉思》

第二种，传统的本体论证明（上帝是完满的，本质中就包含存在Essence includes existence）《第五沉思》

\begin{fancybox}{伽桑狄的反驳}
    存在不是完满中的一项。Existence cannot be said to exist in a thing like a perfection; and if a thing lacks existence, then it is not just imperfect ou lacking perfection, it is nothing at all.
\end{fancybox}

伽桑狄的反驳未必被所有人接受，有两个主要的困难：不存在的东西能不能谈性质？当代的解决方法以Frege为例子，本质清单里面可以任意添加性质，譬如完美的电脑，现实中事实的存在是一个个例，这个个例是不存在的，没有妨碍在清单里面添加性质。（有人提出，想象一个东西完满+不存在，是不可能办到的）

上帝的真诚和理论意义:Discart认为是脱离了私人心灵的第一步。

\begin{itemize}
    \item 欺骗是一种缺陷，所以上帝不会欺骗。
    \item 可以确信算术，几何等数学真理。（所谓的“理智之光”保证在证明了某个东西之后确证其为真，由于上帝不欺骗，我们可以确信背过身去还是真的）
    \item 可以确信自身和外部世界的存在。
\end{itemize}

对于借由上帝完成关于数学真理和外部世界的论证，Discart的这种策略受到了大量质疑。

【生】可以借由缸中之脑的挑战：要是上帝的这种概念也是由缸中之脑构想出来的，这一套论证完全没有受到影响。然而对BIV而言这是一个内蕴的反例。对于上帝不会欺骗的论证疑似为断言，如果完满上帝的观念并不是由上帝endow的

更多的问题：

\begin{itemize}
    \item 上帝的真诚如何兼容错误的存在。
    \item 阿诺德卡氏循环：上帝保证清楚明白的知觉为真。但是对上帝存在的证明预设了清楚明白的知觉为真。（反驳：区分一般知觉和特殊知觉，关于上帝的一部分是特殊的知觉）
\end{itemize}

\subsection{永恒真理}

永恒真理一词在中世纪已经有之，是一个和上帝紧密联系的概念。

\paragraph{对永恒真理概念的改造}

对永恒真理概念的改造；

\begin{itemize}
    \item 永恒真理不同于物质世界或者任何心灵（在这个意义上更加接近于柏拉图主义）
    \item 由上帝创造，依赖于上帝的意志。
    \item 除了逻辑和数学真理之外，包括物理定律
\end{itemize}

为什么永恒真理不能被改写？因为具有如下的理论作用：

\begin{itemize}
    \item 亚里士多德：基于实体形式的解释模型解释物质对象的性质与行为
    \item 笛卡尔：用永恒真理（进而上帝意志）进行解释。
\end{itemize}

\paragraph{卡氏物理学}

广延是物质的本质：

\begin{itemize}
    \item 广延即在空间中延展
    \item 广延存在者具有大小，形状，运动/静止等特性
    \item 广延是物质的本质，因为其他一切特性都可以丢失
    \item 广延以外的声色味等都是思想
\end{itemize}

\subsection{实体二元论}

笛卡尔文本中的实体指的是一个不需要其他东西存在便能存在的东西。

一般认为Discart是一个实体二元论者，他认为有两种存在者：思维存在者（res cogitans）和广延存在者（res extensa），两种实体对应两种本质。

和其他的Dualism相似，这一套二元论也面临mind-body problem。也就是心灵和身体的关系问题，解释心灵如何控制物理实体的运动是这种二元论的重要任务。这是一个小人错误：是“我”这个整体看到了东西，Subject是“我”，不是所谓“心灵”看到了东西。

